% from ega4-8-fr.tex
\oldpage[IV-3]{33}
$Z_{rs}\cap Z_{rt}$~; il résulte aussitôt des définitions que $h'_{sr}$ est un
isomorphisme de $Z_{rs}\cap Z_{rt}$ sur $Z_{sr}\cap Z_{st}$, et que l'on a la
relation $h'_{ts}\circ h'_{sr}=h'_{tr}$. En vertu de (8.8.2, (i)), il existe
donc $\rho\geq\nu$ tel que, pour tout triplet $(r,s,t)$, on ait
$h'_{ts\rho}\circ h'_{sr\rho}=h'_{tr\rho}$. On en conclut que les
isomorphismes $h_{sr\rho}$ vérifient la condition de recollement
$(0_{\mathrm I}, 4.1.7)$ et définissent donc un préschéma $X_\rho$ tel que $X$
soit isomorphe à $X_\rho\times_{S_\rho}S$. En outre, les $Z_{r\rho}$ sont de
présentation finie sur $S_\rho$ et, si on les identifie à des ouverts de
$X_\rho$, les intersections $Z_{r\rho}\cap Z_{s\rho}$, isomorphes aux
$Z_{rs\rho}$, sont quasi-compactes, donc (I, 1.6.3) $X_\rho$ est de
présentation finie sur $S_\rho$.
\end{proof}

\begin{env}[8.8.3]
\label{IV.8.8.3-fr}
Le contenu essentiel de (8.8.2) s'exprime encore en disant que si $S_0$ est
quasi-compact et quasi-séparé, la catégorie des $S$-préschémas de présentation
finie est déterminée à équivalence près par la donnée des catégories des
$S_\lambda$-préschémas de présentation finie, des foncteurs
$\rho_{\mu\lambda}:X_\lambda\rightsquigarrow
X_\lambda\times_{S_\lambda}S_\mu$ (pour $\lambda\leq\mu$) entre ces catégories
et des isomorphismes de transitivité
$\rho_{\nu\mu}\circ\rho_{\mu\lambda}\xrightarrow{\sim}\rho_{\nu\lambda}$. De
façon imagée, on peut dire que la donnée d'un $S$-préschéma de présentation
finie $X$ équivaut «~fonctoriellement~» à la donnée d'un
$S_\lambda$-préschéma de présentation finie $X'_\lambda$~; si un
$S_\mu$-préschéma de présentation finie $X''_\mu$ est tel que $X$ soit
$S$-isomorphe à $X''_\mu\times_{S_\mu}S$, il existe un $\nu$ tel que
$\lambda\leq\nu$, $\mu\leq\nu$ et tel que
$X'_\lambda\times_{S_\lambda}S_\nu$ et
$X''_\mu\times_{S_\mu}S_\nu$ soient $S_\nu$-isomorphes. Le fait que $L$ est
filtrant entraîne en outre que si on se donne une famille \emph{finie}
$(X^{(i)})_{i\in I}$ de $S$-préschémas de présentation finie, et une famille
\emph{finie} $(f^{(j)})_{j\in J}$ de $S$-morphismes entre ces préschémas
($f^{(j)}$ étant donc de la forme
$X^{(\sigma(j))}\to X^{(\tau(j))}$ où $\sigma$ et $\tau$ sont deux
applications de $J$ dans $I$), alors il y a un indice $\lambda\in L$, une
famille $(X_\lambda^{(i)})_{i\in I}$ de $S_\lambda$-préschémas de présentation
finie et une famille $(f_\lambda^{(j)})_{j\in J}$ de $S_\lambda$-morphismes
tels que $X^{(i)}$ s'identifie à
$X_\lambda^{(i)}\times_{S_\lambda}S$ et $f^{(j)}$ à
$f_\lambda^{(j)}\times1_S$ quels que soient $i$ et $j$~; en outre, si on a
une relation $f^{(j)}=f^{(k)}$, on peut supposer $\lambda$ choisi de telle
sorte que $f_\lambda^{(j)}=f_\lambda^{(k)}$.

Considérons en particulier trois $S$-préschémas de présentation finie $X$,
$Y$, $Z$ et deux $S$-morphismes $f:X\to Z$, $g:Y\to Z$, de sorte que le
produit $X\times_ZY$ relatif à ces morphismes est encore un $S$-préschéma de
présentation finie (I, 1.6.2). Alors il résulte de ce qui précède et de
(I, 3.3.11) que l'on peut écrire
$X\times_ZY=(X_\lambda\times_{Z_\lambda}Y_\lambda)\times_{S_\lambda}S$ pour
un $\lambda$ convenable, $X_\lambda$, $Y_\lambda$, $Z_\lambda$ étant des
$S_\lambda$-préschémas de présentation finie~; on peut donc dire que la
détermination des $S$-préschémas de présentation finie par la donnée des
$S_\lambda$-préschémas de présentation finie est «~\emph{compatible avec les
produits fibrés}~». On a vu d'autre part (8.6.3) que si $g$ est une immersion,
on peut supposer qu'il en est de même de $g_\lambda:Y_\lambda\to Z_\lambda$~;
identifiant $Y$ (resp. $Y_\lambda$) avec un sous-préschéma de $Z$ (resp.
$Z_\lambda$), on voit donc que l'on peut écrire, pour un $\lambda$ convenable,
$f^{-1}(Y)=f_\lambda^{-1}(Y_\lambda)\times_{S_\lambda}S$ (I, 4.4.1)~; il y a
donc aussi «~\emph{compatibilité avec la formation des images réciproques de
sous-préschémas}~». Plus particulièrement, si $f_1$, $f_2$ sont deux
$S$-morphismes de $X$ dans $Y$, on appelle \emph{noyau} de ce couple de
morphismes l'image réciproque $N$ du sous-préschéma diagonal de $Y\times_SY$
par le $S$-morphisme $(f_1,f_2)_S$~; on déduit de ce qui précède que l'on aura
pour $\lambda$ convenable $N=N_\lambda\times_{S_\lambda}S$, où $N_\lambda$
est le noyau du couple de $S_\lambda$-morphismes
$(f_{1\lambda},f_{2\lambda})$ de $X_\lambda$ dans $Y_\lambda$. Ces remarques
s'étendent à des produits finis quelconques et aux «~noyaux~» de systèmes
finis quelconques de $S$-morphismes d'un $S$-préschéma de présentation finie
